\section{Introduction}
\label{sec:intro}

Agricultural machinery has become a central carrier of field-scale production data. Modern combine harvesters continuously record GNSS trajectories, visual observations, and machine sounds while executing harvesting, transfer, waiting, unloading, and road-driving operations. Turning these raw sensor streams into reliable operational labels is the core goal of \task{}. Accurate labels support operation-efficiency evaluation, harvest-loss analysis, machine scheduling, and downstream digital-farm management.

Most prior work treats agricultural machinery trajectory analysis as field-road classification or coarse-grained operational state classification~\cite{pan2024chaos,zhai2024atrnet,zhai2024bilstm}. These settings are important, but they do not capture the operational detail required by real combine-harvester workflows. In this work, we study an eleven-class problem that distinguishes reverse empty harvesting, straight empty harvesting, turning empty harvesting, full-load harvesting, reverse transfer, straight transfer, turning transfer, engine-off waiting, idling waiting, unloading, and road driving.

The main challenge is that several classes are not separable from trajectory geometry alone. Engine-off waiting, idling waiting, and unloading are all nearly stationary in GNSS coordinates. Straight empty harvesting and straight transfer can both appear as smooth linear motion. Turning empty harvesting and turning transfer can have similar curvature. In these cases, the decisive evidence lies outside the trajectory stream: video reveals crop rows, header state, and scene context, while audio reflects engine load, idling, unloading mechanisms, and other process-level cues. This motivates a multimodal formulation of \task{}.

This paper is written around the current controlled multimodal eleven-class classification suite with matched temporal context. The trajectory modality uses a 128-second causal sequence, selected because a completed seq\_len ablation found it to provide the best mean macro-F1 among 64, 128, 256, and 512 seconds. The ViT visual modality follows the same 128-second causal window and selects the nine frames closest to the label time. The audio modality uses AST-based spectrogram encoding. The current \ours{} model replaces direct feature concatenation with class-adaptive logit-level fusion, so that different operation classes can assign different weights to BiLSTM, ViT, and AST evidence.

The contributions of this paper are:
\begin{enumerate}
    \item We formulate an eleven-class \task{} benchmark for combine-harvester operations and define a unified BiLSTM--ViT--AST experimental protocol with matched causal temporal context.
    \item We evaluate modality-specific baselines, including BiLSTM trajectory encoding, pretrained ViT visual encoding with temporal pooling, and AST audio encoding, under a shared split and fixed seed.
    \item We propose a class-adaptive trimodal fusion model that initializes from unimodal checkpoints and combines modality-specific logits with class-wise BiLSTM, ViT, and AST gates, improving fixed-seed macro-F1 by 2.44 percentage points over direct trimodal concatenation.
    \item We provide trajectory-level diagnostic visualizations, per-class confusion analysis, fusion ablations, and sequence-length ablations showing that multimodal fusion improves most harvesting, transfer, and waiting/transfer classes while revealing unloading and road-driving states that remain challenging.
\end{enumerate}

The remainder of this paper is organized as follows. \Cref{sec:related} reviews agricultural machinery trajectory analysis and multimodal learning. \Cref{sec:method} describes the proposed multimodal model. \Cref{sec:exp} presents the dataset, implementation details, and experiments. \Cref{sec:discussion} discusses modality-specific findings and limitations. \Cref{sec:conclusion} concludes the paper.
